% Ad Familiares 12.15 (ad-familiares-12-15)
% Source: https://raw.githubusercontent.com/PerseusDL/canonical-latinLit/master/data/phi0474/phi056/phi0474.phi056.perseus-lat1.xml
% Fetched by scripts/fetch_latin.py

% Dateline (Perseus): Scr. Pergae 1-6 §§ iiii K., § 7 iiii Non. Iun. a. 711 (43).

\ciceroLetterOpener{P. LENTVLVS P. F. PROQ. PROPR. S. D. COSS. PR. TR. PL. SENATVI POPVLO PLEBIQVE ROMANAE}

\ciceroSection{1} S. v. l. v. v. b. e. v. scelere Dolabellae oppressa Asia in proximam provinciam Macedoniam praesidiaque rei p. quae M. Brutus, v. c., tenebat me contuli et id egi ut, per quos celerrime possent, Asia provincia vectigaliaque in vestram potestatem redigerentur. quod cum pertimuisset Dolabella vastata provincia, correptis vectigalibus, praecipue civibus Romanis omnibus crudelissime denudatis ac divenditis celeriusque Asia excessisset quam eo praesidium adduci potuisset, diutius morari aut exspectare praesidium non necesse habui et quam primum ad meum officium revertendum mihi esse existimavi, ut et reliqua vectigalia exigerem et quam deposui pecuniam conligerem, quicquid ex ea correptum esset aut quorum id culpa accidisset cognoscerem quam primum et vos de omni re facerem certiores.

\ciceroSection{2} interim cum per insulas in Asiam naviganti mihi nuntiatum esset classem Dolabellae in Lycia esse Rhodiosque navis compluris instructas et paratas in aqua habere, cum Hs navibus quas aut mecum adduxeram aut comparaverat Patiscus proq., homo mihi cum familiaritate tum etiam sensibus in re p. coniunctissimus, Rhodum deverti confisus auctoritate vestra senatusque consulto quo hostem Dolabellam iudicaratis, foedere quoque quod cum iis M. Marcello, Ser. Sulpicio coss. renovatum erat, quo iuraverant Rhodii eosdem hostis se habituros quos senatus populusque R. quae res nos vehementer fefellit ; tantum enim afuit ut illorum praesidio nostram firmaremus classem, ut etiam a Rhodiis urbe, portu, statione quae extra urbem est, commeatu, aqua denique prohiberentur nostri milites, nos vix ipsi singulis cum navigiolis reciperemur. quam indignitatem deminutionemque t maiestatis non solum iuris nostri sed etiam imperi populique Romani idcirco tulimus quod interceptis litteris cognoramus Dolabellam, si desperasset de Syria Aegyptoque, quod necesse erat fieri, in navis cum omnibus suis latronibus atque omni pecunia conscendere esse paratum Italiamque petere ; idcirco etiam navis onerarias, quarum minor nulla erat duum milium amphorum, contractas in Lycia a classe eius obsideri.

\ciceroSection{3} huius rei timore, p. c., percitus iniurias perpeti et cum contumelia etiam nostra omnia prius experiri malui. itaque ad illorum voluntatem introductus in urbem et in senatum eorum quam diligentissime potui causam rei p. egi periculumque omne quod instaret si ille latro cum suis omnibus navis conscendisset exposui. Rhodios autem tanta in pravitate animadverti ut omnis firmiores putarent quam bonos, ut hanc concordiam et conspirationem omnium ordinum ad defendendam libertatem propense non crederent esse factam, ut patientiam senatus et optimi cuiusque manere etiam nunc confiderent nec potuisse audere quemquam Dolabellam hostem iudicare, ut denique omnia quae improbi fingebant magis vera existimarent quam quae vere facta erant et a nobis docebantur.

\ciceroSection{4} qua mente etiam ante nostrum adventum post Treboni indignissimam caedem ceteraque tot tamque nefaria facinora binae profectae erant ad Dolabellam legationes eorum, et quidem novo exemplo, contra leges ipsorum, prohibentibus iis qui tum magistratus gerebant. † haec sive timore, ut dictitant, de agris quos in continenti habent sive furore sive patientia paucorum, qui et antea pari contumelia viros clarissimos adfecerant et nunc maximos magistratus gerentes nullo exemplo neque nostra ex parte neque nostro praesentium neque imminenti Italiae urbique nostrae periculo, si ille parricida cum suis latronibus navibus ex Asia Syriaque expulsus Italiam petisset, mederi, cum facile possent, noluerunt†.

\ciceroSection{5} non nullis etiam ipsi magistratus veniebant in suspicionem detinuisse nos et demorati esse, dum classis Dolabellae certior fieret de adventu nostro. quam suspicionem consecutae res aliquot auxerunt, maxime quod subito ex Lycia Sex. Marius et C. Titius, legati Dolabellae, a classe discesserunt navique longa profugerunt onerariis relictis, in quibus conligendis non minimum temporis laborisque consumpserant. itaque cum ab Rhodo cum iis quas habueramus navibus in Lyciam venissemus, navis onerarias recepimus dominisque restituimus idemque, quod maxime verebamur, ne posset Dolabella cum suis latronibus in Italiam venire, timere desumus ; classem fugientem persecuti suiuus usque Sidam, quae extrema regio est provinciae meae.

\ciceroSection{6} ibi cognovi partem navium Dolabellae diffugisse, reliquas Syriam Cyprumque petisse. quibus disiectis, cum scirem C. Cassi, singularis civis et ducis, classem maximam fore praesto in Syria, ad meum officium reverti daboque operam ut meum studium, diligentiam vobis, p. c., reique p. praestem pecuniamque quam maximam potero et quam celerrime cogam omnibusque rationibus ad vos mittam. si percurrero provinciam et cognovero qui nobis et rei p. fidem praestiterunt in conservanda pecunia a me deposita quique scelere ultro deferentes pecuniam publicam hoc munere societatem facinorum cum Dolabella inierunt, faciam vos certiores. de quibus, si vobis videbitur, si, ut meriti sunt, graviter constitueritis nosque vestra auctoritate firmaveritis, facilius et reliqua exigere vectigalia et exacta servare poterimus. interea quo commodius vectigalia tueri provinciamque ab iniuria defendere possim, praesidium voluntarium necessariumque comparavi.

\ciceroSection{7} his litteris scriptis milites circiter xxx, quos Dolabella ex Asia conscripserat, ex Syria fugientes in Pamphyliam venerunt. hi nuntiaverunt Dolabellam Antiocheam quae in Syria est venisse, non receptum conatum esse aliquotiens vi introire ; repulsum semper esse cum magno suo detrimento itaque DC circiter amissis, aegris relictis noctu Antiochea profugisse Laudiceam versus ; ea nocte omnis fere Asiaticos milites ab eo discessisse ; ex his ad octingentos Antiocheam redisse et se iis tradidisse qui a Cassio relicti urbi illi praeerant, ceteros per Amanum in Ciliciam descendisse, quo ex numero se quoque esse dicebant; Cassium autem cum suis omnibus copiis nuntiatum esse quadridui iter Laudicea afuisse tum cum Dolabella eo tenderet. quam ob rem opinione celerius confido sceleratissimum latronem poenas daturum. iiii N. Iun. Pergae.
